% =========================================================
% 06_systems_of_linear_equations.tex
% Vectors and Linear Algebra
% =========================================================

% =========================================================
% SECTION DIVIDER
% =========================================================

\section{Systems of Linear Equations}

% =========================================================
% SYSTEM OF LINEAR EQUATIONS
% =========================================================

\begin{frame}{System of Linear Equations}

A \textbf{linear equation} in variables
$x_1,x_2,\ldots,x_n$ has the form

\[
a_1x_1+a_2x_2+\cdots+a_nx_n=b.
\]

A \textbf{system of linear equations} consists of
multiple linear equations considered simultaneously.

\vspace{0.4cm}

For example,

\[
\begin{cases}
2x+y=5,\\
x-y=1.
\end{cases}
\]

\vspace{0.3cm}

The objective is to find values of the variables that
satisfy \textbf{all equations simultaneously}.

\end{frame}


% =========================================================
% SOLUTION OF A SYSTEM
% =========================================================

\begin{frame}{Solution of a System}

For the system
\vspace{-0.1cm}
\[
\begin{cases}
2x+y=5,\\
x-y=1,
\end{cases}
\]

adding the two equations gives
\vspace{-0.1cm}
\[
3x=6.
\]

Therefore,
\vspace{-0.1cm}
\[
x=2.
\]

Substituting into
\vspace{-0.1cm}
\[
x-y=1
\]

gives
\vspace{-0.1cm}
\[
y=1.
\]

Hence,
\vspace{-0.1cm}
\[
\boxed{
(x,y)=(2,1)
}
\]
\vspace{-0.1cm}
is the solution.

\end{frame}


% =========================================================
% GEOMETRIC INTERPRETATION
% =========================================================

\begin{frame}{Geometric Interpretation}

In two variables, each linear equation represents a line.

For example,

\[
2x+y=5
\]

and

\[
x-y=1
\]

represent two lines in the plane.

\vspace{0.4cm}

Their solution is the point where the two lines intersect.

\[
\boxed{
\text{Solution of system}
=
\text{Common intersection}
}
\]

This geometric interpretation becomes more general when
we move to higher dimensions.

\end{frame}


% =========================================================
% CONSISTENCY
% =========================================================

\begin{frame}{Consistency of a System}

A system of linear equations is \textbf{consistent} if it has
at least one solution.

It is \textbf{inconsistent} if it has no solution.

\vspace{0.5cm}

A consistent system may have:

\[
\boxed{
\text{Exactly one solution}
}
\]

or

\[
\boxed{
\text{Infinitely many solutions}.
}
\]

\vspace{0.4cm}

Therefore, the three possibilities are:

\[
\boxed{
\text{Unique solution}
\qquad
\text{Infinitely many solutions}
\qquad
\text{No solution}
}
\]

\end{frame}


% =========================================================
% CONSISTENCY GEOMETRICALLY
% =========================================================

\begin{frame}{Consistency: Geometric View}

For two equations in two variables:

\vspace{0.3cm}

\begin{columns}[T,totalwidth=\textwidth]

\begin{column}{0.30\textwidth}

\begin{block}{One Solution}
Two distinct lines intersect at one point.
\end{block}

\end{column}

\begin{column}{0.30\textwidth}

\begin{block}{No Solution}
Parallel distinct lines never intersect.
\end{block}

\end{column}

\begin{column}{0.30\textwidth}

\begin{block}{Infinite Solutions}
Coincident lines represent the same equation.
\end{block}

\end{column}

\end{columns}

\vspace{0.5cm}

Thus, consistency is closely related to the geometry of the
equations.

\end{frame}


% =========================================================
% MATRIX FORM
% =========================================================

\begin{frame}{Matrix Form of a Linear System}

Consider

\[
\begin{cases}
a_{11}x_1+a_{12}x_2=b_1,\\
a_{21}x_1+a_{22}x_2=b_2.
\end{cases}
\]

Define

\[
A=
\begin{bmatrix}
a_{11}&a_{12}\\
a_{21}&a_{22}
\end{bmatrix},
\qquad
\vect{x}
=
\begin{bmatrix}
x_1\\
x_2
\end{bmatrix},
\qquad
\vect{b}
=
\begin{bmatrix}
b_1\\
b_2
\end{bmatrix}.
\]

Then the system becomes

\[
\boxed{
A\vect{x}=\vect{b}.
}
\]

\end{frame}


% =========================================================
% Ax = b
% =========================================================

\begin{frame}{The Equation $A\vect{x}=\vect{b}$}

The matrix equation

\[
\boxed{
A\vect{x}=\vect{b}
}
\]

can be interpreted as a transformation.

\[
\vect{x}
\xrightarrow{\quad A\quad}
A\vect{x}
=
\vect{b}.
\]

\vspace{0.5cm}

The problem is therefore:

\[
\boxed{
\text{Find the input }\vect{x}
\text{ that produces the output }\vect{b}.
}
\]

\vspace{0.4cm}

This connects systems of equations directly with
\textbf{linear transformations}.

\end{frame}


% =========================================================
% COLUMN INTERPRETATION
% =========================================================

\begin{frame}{$A\vect{x}=\vect{b}$ as a Linear Combination}

Let
\vspace{-1.9cm}
\begin{columns}[c,totalwidth=\textwidth]

% ---------- LEFT ----------
\begin{column}{0.45\textwidth}

\begin{center}

\[
A=
\begin{bmatrix}
|&|& &|\\
\vect{a}_1&\vect{a}_2&\cdots&\vect{a}_n\\
|&|& &|
\end{bmatrix}
\],

\end{center}

\end{column}

% ---------- RIGHT ----------
\begin{column}{0.45\textwidth}

\begin{center}

\[
\vect{x}
=
\begin{bmatrix}
x_1\\
x_2\\
\vdots\\
x_n
\end{bmatrix}
\]

\end{center}

\end{column}

\end{columns}

\vspace{0.1cm}

Then
\vspace{-0.1cm}
\[
A\vect{x}
=
x_1\vect{a}_1+
x_2\vect{a}_2+
\cdots+
x_n\vect{a}_n.
\]

Therefore,

\[
A\vect{x}=\vect{b}
\]

asks whether $\vect{b}$ can be expressed as a linear
combination of the columns of $A$.

\vspace{-0.1cm}

\begin{center}

\[
\boxed{
\vect{b}\in\operatorname{Col}(A)
}
\]

\end{center}

This is precisely the condition for consistency.

\end{frame}


% =========================================================
% AUGMENTED MATRIX
% =========================================================

\begin{frame}{Augmented Matrix}

A system can be represented using an \textbf{augmented matrix}.

For

\[
\begin{cases}
2x+y=5,\\
x-y=1,
\end{cases}
\]

the augmented matrix is

\[
\boxed{
\left[
\begin{array}{cc|c}
2&1&5\\
1&-1&1
\end{array}
\right].
}
\]

\vspace{0.4cm}

The vertical line separates:

\[
\text{coefficient matrix}
\qquad\text{from}\qquad
\text{right-hand side}.
\]

\end{frame}


% =========================================================
% ELEMENTARY ROW OPERATIONS
% =========================================================

\begin{frame}{Elementary Row Operations}

Gaussian elimination uses three elementary row operations.

\vspace{0.4cm}

\begin{enumerate}
    \item Interchange two rows.
    \item Multiply a row by a non-zero scalar.
    \item Add a multiple of one row to another row.
\end{enumerate}

\vspace{0.5cm}

These operations transform the system into an equivalent
system with the same solution set.

\[
\boxed{
\text{Equivalent system}
\quad\Longrightarrow\quad
\text{Same solutions}
}
\]

\end{frame}


% =========================================================
% GAUSSIAN ELIMINATION
% =========================================================

\begin{frame}{Gaussian Elimination}

Gaussian elimination transforms an augmented matrix into
\textbf{row echelon form}.

Consider
\vspace{-0.1cm}
\[
\left[
\begin{array}{cc|c}
2&1&5\\
1&-1&1
\end{array}
\right].
\]

Swap the rows:
\vspace{-0.1cm}
\[
\left[
\begin{array}{cc|c}
1&-1&1\\
2&1&5
\end{array}
\right].
\]

Then perform

\[
R_2\leftarrow R_2-2R_1.
\]

This gives

\[
\left[
\begin{array}{cc|c}
1&-1&1\\
0&3&3
\end{array}
\right].
\]

\end{frame}


% =========================================================
% BACK SUBSTITUTION
% =========================================================

\begin{frame}{Gaussian Elimination: Back Substitution}

From
\vspace{-0.1cm}
\[
\left[
\begin{array}{cc|c}
1&-1&1\\
0&3&3
\end{array}
\right]
\]

we obtain
\vspace{-0.1cm}
\[
3y=3
\]

and therefore
\vspace{-0.1cm}
\[
y=1.
\]

Substituting into the first equation,
\vspace{-0.2cm}
\[
x-y=1
\]

gives
\vspace{-0.2cm}
\[
x=2.
\]

Hence,
\vspace{-0.2cm}
\[
\boxed{
\vect{x}
=
\begin{bmatrix}
2\\
1
\end{bmatrix}.
}
\]

\end{frame}


% =========================================================
% ROW ECHELON FORM
% =========================================================

\begin{frame}{Row Echelon Form}

A matrix is in \textbf{row echelon form} when:

\begin{itemize}
    \item all non-zero rows occur above zero rows,
    \item each leading entry is to the right of the leading
    entry in the row above it,
    \item entries below each leading entry are zero.
\end{itemize}

For example,

\[
\begin{bmatrix}
1&2&-1\\
0&1&3\\
0&0&1
\end{bmatrix}
\]

is in row echelon form.

\vspace{0.4cm}

Gaussian elimination systematically produces this structure.

\end{frame}


% =========================================================
% REDUCED ROW ECHELON FORM
% =========================================================

\begin{frame}{Reduced Row Echelon Form}

A matrix is in \textbf{reduced row echelon form} when, in addition
to row echelon conditions:

\begin{itemize}
    \item every leading entry is $1$,
    \item each leading $1$ is the only non-zero entry in its column.
\end{itemize}

For example,

\[
\boxed{
\begin{bmatrix}
1&0&2\\
0&1&-3\\
0&0&0
\end{bmatrix}
}
\]

is in reduced row echelon form.

\vspace{0.3cm}

RREF makes the solution structure of a system easier to identify.

\end{frame}


% =========================================================
% UNIQUE SOLUTION
% =========================================================

\begin{frame}{Identifying a Unique Solution}

Consider

\[
\left[
\begin{array}{ccc|c}
1&0&2&4\\
0&1&-1&3\\
0&0&1&2
\end{array}
\right].
\]

Every variable has a pivot.

\vspace{0.4cm}

Therefore, the system has

\[
\boxed{
\text{exactly one solution}.
}
\]

In general, a unique solution occurs when every variable
is determined by the equations.

\end{frame}


% =========================================================
% NO SOLUTION
% =========================================================

\begin{frame}{Detecting No Solution}

Suppose elimination produces
\vspace{-0.1cm}
\[
\left[
\begin{array}{ccc|c}
1&2&0&3\\
0&0&0&5
\end{array}
\right].
\]

The second row represents
\vspace{-0.1cm}
\[
0=5.
\]

This is impossible.

Therefore,
\vspace{-0.1cm}
\[
\boxed{
\text{The system is inconsistent.}
}
\]

A row of the form

\[
\left[
\begin{array}{ccc|c}
0&0&0&c
\end{array}
\right],
\qquad c\neq0,
\]

signals that the system has \textbf{no solution}.

\end{frame}


% =========================================================
% INFINITELY MANY SOLUTIONS
% =========================================================

\begin{frame}{Infinitely Many Solutions}

Consider a system whose reduced form contains
\vspace{-0.1cm}
\[
\left[
\begin{array}{ccc|c}
1&0&2&4\\
0&1&-1&3\\
0&0&0&0
\end{array}
\right].
\]

The third row gives no new constraint. One variable can therefore be chosen freely.

Let
\vspace{-0.2cm}
\[
z=t.
\]

Then
\vspace{-0.1cm}
\[
x=4-2t,
\qquad
y=3+t.
\]

Hence,
\vspace{-0.1cm}
\[
\boxed{
\begin{bmatrix}
x\\y\\z
\end{bmatrix}
=
\begin{bmatrix}
4\\3\\0
\end{bmatrix}
+
t
\begin{bmatrix}
-2\\1\\1
\end{bmatrix}.
}
\]
\vspace{-0.1cm}
There are infinitely many solutions.

\end{frame}


% =========================================================
% INVERSE MATRIX
% =========================================================

\begin{frame}{Solving $A\vect{x}=\vect{b}$ Using an Inverse}
\vspace{-0.1cm}
Suppose $A$ is a square matrix with an inverse.

Starting with

\[
A\vect{x}=\vect{b},
\]

multiply both sides by $A^{-1}$:

\[
A^{-1}A\vect{x}
=
A^{-1}\vect{b}.
\]

Since
\vspace{-0.1cm}
\[
A^{-1}A=I,
\]

we obtain
\vspace{-0.1cm}
\[
\boxed{
\vect{x}=A^{-1}\vect{b}.
}
\]

This gives a direct formula for the solution.

\end{frame}


% =========================================================
% CONDITIONS FOR INVERSE SOLUTION
% =========================================================

\begin{frame}{When Does $A^{-1}\vect{b}$ Exist?}

The expression

\[
\vect{x}=A^{-1}\vect{b}
\]

requires $A$ to be:

\begin{itemize}
    \item square,
    \item invertible.
\end{itemize}

For a square matrix,

\[
A^{-1}\text{ exists}
\quad\Longleftrightarrow\quad
\det(A)\neq0.
\]


When these conditions hold, the system

\[
A\vect{x}=\vect{b}
\]

has a unique solution for every $\vect{b}$.

\end{frame}


% =========================================================
% GAUSSIAN ELIMINATION VS INVERSE
% =========================================================

\begin{frame}{Two Ways to Solve a Linear System}

For a system

\[
A\vect{x}=\vect{b},
\]

two approaches are:

\vspace{0.3cm}

\begin{columns}[T,totalwidth=\textwidth]

\begin{column}{0.44\textwidth}

\begin{block}{Gaussian Elimination}

Transform the augmented matrix and solve the resulting system.

\[
[A\mid\vect{b}]
\]

\end{block}

\end{column}

\begin{column}{0.44\textwidth}

\begin{block}{Inverse Method}

If $A$ is invertible,

\[
\vect{x}=A^{-1}\vect{b}.
\]

\end{block}

\end{column}

\end{columns}

\vspace{0.5cm}

For computational problems involving many right-hand sides,
factorization-based methods are generally more useful than
explicitly computing $A^{-1}$.

\end{frame}


% =========================================================
% HOMOGENEOUS SYSTEM
% =========================================================

\begin{frame}{Homogeneous Linear Systems}

A homogeneous system has the form

\[
\boxed{
A\vect{x}=\mathbf{0}.
}
\]

It always has at least the trivial solution

\[
\boxed{
\vect{x}=\mathbf{0}.
}
\]

\vspace{0.4cm}

Depending on $A$, there may also be non-zero solutions.

\vspace{0.4cm}

\[
\boxed{
\text{Trivial solution}
\qquad\text{vs.}\qquad
\text{Non-trivial solutions}
}
\]

This will become important when studying the
\textbf{null space} and linear dependence.

\end{frame}